% Corrigé intégré automatiquement pour la banque Exercices.
% Source : A_Integrer/DDS_Reserve/982_TEC/982_TEC.tex
\begin{corrigebox}[Corrigé]
\CorrigeQuestion{1}
$I_{eq}(m)=I_S+I_R\left(N+1\right)^2+M_y\left(N_1\left(\theta_m\right)\right)^2$, $A=N$, $B=\dfrac{M_y N_1\left( \theta_m\right)N'_1\left( \theta_m\right)}{N}$ et $C\left( \theta_m\right)=-N_1\left( \theta_m\right)$.
\CorrigeQuestion{2}
$AC_{m0}-B\left( \theta_{m0}\right) \Omega_{m0}^2-C\left( \theta_{m0}\right)F_0 = 0 $ et $C_{m0}(t)=\dfrac{K_C}{R}\left(U_0-K_e\Omega_{m0}\right)$
\CorrigeQuestion{3}
$\dfrac{RI_{eq}\left( \theta_m\right)}{N\left( AK_cK_e+2RB\left( \theta_m\right)\Omega_{m0}\right)} \dot{\Omega}_1+\Omega_1 $

 $= 
\dfrac{AK_c}{ AK_cK_e+2RB\left( \theta_m\right)\Omega_{m0}} u_1
- \dfrac{RC\left( \theta_m\right)}{ AK_cK_e+2RB\left( \theta_m\right)\Omega_{m0}} f_1$
\CorrigeQuestion{4}
$\dfrac{\Omega_1(p)}{U_1(p)}=\dfrac{K_0}{1+\tau_0 p}$ et $\dfrac{\Omega_1(p)}{F_1(p)}=-\dfrac{D_0}{1+\tau_0 p}$.
\CorrigeQuestion{5}
On réduit le schéma-blocs de l'intérieur vers l'extérieur : associations en série par produit, associations en parallèle par somme et boucle de retour par $H_{BF}=\dfrac{G}{1+GH}$ (retour négatif). Après simplification, on ordonne le numérateur et le dénominateur selon les puissances de $p$, puis on identifie le gain statique, la classe et les paramètres canoniques.
\end{corrigebox}
